\section{Forme de Weil, double cercle et représentation spectrale}
\label{lectura:manuscrito-03}
La factorisation et les horloges de la section~\ref{lectura:manuscrito-02} produisent deux lectures d’une même arithmétique. La première décompose une forme normale en irréductibles ; la seconde enregistre les retours de leurs actions multiplicatives. Le développement spectral doit conserver les deux provenances. Son objet ne s’obtient pas en assignant à un retour un zéro connu d’une fonction : on construit les fonctionnelles, on établit leur compatibilité et on identifie ensuite leur réalisation analytique.

Ce chapitre réunit les développements conservés de la forme complétée et du double cercle. Il expose séparément les résultats inconditionnels de normalisation, les conséquences exactes d'une identité positive de frontière et les réalisations de cette identité. Cette séparation fixe l'ordre logique des démonstrations. En particulier, le calcul d'une représentation utilisant la positivité ne peut remplacer la construction qui doit l'établir.

\subsection{Du raffinement cyclique à la forme complétée}
Soit \(N=3^m\), avec \(m\geq1\). Sur le cycle de \(N\) positions, on conserve l'opérateur quadratique du chapitre précédent. Si les indices de Fourier sont choisis dans l'intervalle symétrique

\[
-\frac{N-1}{2}\leq n\leq\frac{N-1}{2},
\]
ses valeurs propres quadratiques sont

\[
\lambda_{m,n}=\frac{N^2}{\pi}\sin^2\!\left(\frac{\pi n}{N}\right).
\]
La coordonnée \(\pi\) est la sortie régionale antérieure du noyau HMT. À cette étape, elle intervient dans la normalisation de la réalisation analytique, et non dans la sélection des irréductibles. Le mode constant est retiré et les deux orientations d'un mode non nul ne sont comptées qu'une seule fois :

\[
\Theta_m(x)=\frac12\sum_{n\neq0}e^{-x\lambda_{m,n}}
=\sum_{n=1}^{(N-1)/2}e^{-x\lambda_{m,n}},\qquad x>0.
\]
\HMTresultado{Proposition}{Limite de la trace de chaleur}{rz-num-03-1}
Pour tout \(x>0\),

\[
\lim_{m\to\infty}\Theta_m(x)
=\Theta_+(x):=\sum_{n=1}^{\infty}e^{-\pi n^2x}.
\]
La convergence est uniforme lorsque \(x\) parcourt un compact de \((0,\infty)\).

\textbf{Démonstration.} Pour un indice fixé \(n\), l'équivalence \(\sin y/y\to1\) donne \(\lambda_{m,n}\to\pi n^2\). De plus, \(\sin y\geq2y/\pi\) pour \(0\leq y\leq\pi/2\), d'où

\[
\lambda_{m,n}\geq\frac{4n^2}{\pi}.
\]
Si \(x\geq a>0\), chaque terme est dominé par \(e^{-4an^2/\pi}\), dont la somme converge. En prolongeant par zéro les termes hors de l’intervalle d’indices du cycle et en appliquant la convergence dominée, on obtient la limite et son uniformité sur les compacts. \hfill\(\square\)

L'inversion de cette trace explique pourquoi la complétion contient une contribution polaire en plus des modes positifs. Définissons

\[
\vartheta(x)=1+2\Theta_+(x).
\]
La périodisation de la gaussienne

\[
F_x(t)=\sum_{n\in\mathbb Z}e^{-\pi x(t+n)^2}
\]
converge avec toutes ses dérivées. Son coefficient de Fourier d'indice \(k\) est \(x^{-1/2}e^{-\pi k^2/x}\), obtenu en intégrant la gaussienne sur la droite. La série de Fourier converge absolument ; en l'évaluant en \(t=0\), on obtient

\[
\vartheta(x)=x^{-1/2}\vartheta(1/x).
\]
En termes de trace positive,

\[
\Theta_+(x)=x^{-1/2}\Theta_+(1/x)+\frac{x^{-1/2}-1}{2}.
\]
Le dernier terme conserve la différence entre retirer le mode constant avant ou après avoir inversé l’échelle.

\HMTresultado{Proposition}{Transformée de Mellin et prolongement}{rz-num-03-2}
Pour \(\Re s>1\), la fonction de partition arithmétique \(Z(s)=\sum_{n\geq1}n^{-s}\) satisfait

\[
\Lambda(s):=\int_0^\infty\Theta_+(x)x^{s/2-1}\,dx
=\pi^{-s/2}\Gamma(s/2)Z(s).
\]
La fonction

\[
\xi(s)=\frac12s(s-1)\Lambda(s)
\]
admet un prolongement entier et satisfait \(\xi(s)=\xi(1-s)\).

\textbf{Démonstration.} La convergence absolue pour \(\Re s>1\) permet d'intégrer terme par terme. Le changement \(u=\pi n^2x\) donne le facteur \(\pi^{-s/2}\Gamma(s/2)n^{-s}\). Pour prolonger l'intégrale, on sépare les intervalles \((0,1)\) et \((1,\infty)\) et on applique l'inversion précédente au premier. Il en résulte

\[
\Lambda(s)=\frac1{s(s-1)}
+\int_1^\infty\Theta_+(x)
\left(x^{s/2-1}+x^{(1-s)/2-1}\right)\,dx.
\]
L’intégrale est entière en \(s\) : la décroissance exponentielle de \(\Theta_+(x)\) domine uniformément toute puissance et toute puissance de \(\log x\) sur les compacts du plan. Le facteur \(s(s-1)\) élimine les deux pôles montrés explicitement. L’expression est invariante sous \(s\mapsto1-s\), ce qui prouve l’équation fonctionnelle. \hfill\(\square\)

La transformée de Mellin de la trace n'est pas simplement \(Z(s)\) : elle contient les facteurs \(\pi^{-s/2}\Gamma(s/2)\). Il est indispensable de les maintenir explicites, car leurs contributions réapparaissent dans la forme quadratique complétée. Le produit sur les irréductibles du chapitre précédent et cette construction additive coïncident dans le demi-plan de convergence par une identité démontrée de fonctionnelles.

\subsection{Les contributions arithmétique, gamma et polaire}
\label{rz:forma}
Fixons \(R>0\) et considérons

\[
\mathcal D_R=C_c^\infty((-R/2,R/2)).
\]
Le produit scalaire sera linéaire en la première variable. On utilise

\[
\widehat f(z)=\int_{\mathbb R}f(x)e^{-izx}\,dx,\qquad
(T_af)(x)=f(x-a),
\]
\[
C_{f,g}(a)=\langle f,T_ag\rangle_{L^2(\mathbb R)},\qquad
h_{f,g}(z)=\widehat f(z)\overline{\widehat g(\overline z)}.
\]
Ces définitions fixent simultanément l'orientation des translations et la normalisation de Fourier. Pour \(f,g\in\mathcal D_R\), \(C_{f,g}(a)=0\) si \(|a|\geq R\).

Les irréductibles déjà construits fournissent les indices premier–puissance et leurs poids :

\[
\ell_{p,k}=k\log p,\qquad
w_{p,k}=(\log p)p^{-k/2},\qquad k\geq1.
\]
La contribution provenant de la complétion gamma utilise

\[
\mu(a)=\frac{e^{-a/2}}{1-e^{-2a}},\qquad
c_\Gamma=\psi(1/4)-\log\pi,
\]
où \(\psi=\Gamma'/\Gamma\). En particulier, \(c_\Gamma<0\). Par exemple, la monotonie de \(\psi\) sur l’axe positif, obtenue à partir de \(\psi'(x)=\sum_{n\geq0}(n+x)^{-2}>0\), et \(\psi(1)=-\gamma_{\mathrm E}<0\) donnent cette inégalité.

La forme complétée est

\[
\begin{aligned}
\mathscr W_R(f,g)={}&h_{f,g}(i/2)+h_{f,g}(-i/2)
+c_\Gamma C_{f,g}(0)\\
&+\int_0^\infty\mu(a)
\bigl(2C_{f,g}(0)-C_{f,g}(a)-C_{f,g}(-a)\bigr)\,da\\
&-\sum_{\ell_{p,k}\leq R}w_{p,k}
\bigl(C_{f,g}(\ell_{p,k})+C_{f,g}(-\ell_{p,k})\bigr).
\end{aligned}
\]
La somme premier–puissance est finie : seuls interviennent \(p^k\leq e^R\). La forme n'a pas été remplacée par un échantillonnage fini de zéros.

\HMTresultado{Proposition}{Existence et compatibilité de la forme}{rz-num-03-3}
L'expression précédente définit une forme sesquilinéaire hermitienne et continue sur \(\mathcal D_R\). Si \(R<R'\), ses restrictions coïncident :

\[
\mathscr W_{R'}|_{\mathcal D_R\times\mathcal D_R}=\mathscr W_R.
\]
\textbf{Démonstration.} Puisque les translations sont unitaires,

\[
2C_{f,g}(0)-C_{f,g}(a)-C_{f,g}(-a)
=\langle(I-T_a)f,(I-T_a)g\rangle.
\]
Pour \(a\to0\), chaque différence a pour norme \(O(a)\), tandis que \(\mu(a)=O(a^{-1})\). L’intégrande est donc \(O(a)\). Pour \(a\to\infty\), les différences ont des normes bornées et \(\mu(a)=O(e^{-a/2})\). L’intégrale converge absolument. Ces mêmes estimations, exprimées au moyen des normes de \(f,g\) et de leurs premières dérivées sur un compact fixé, prouvent la continuité. La symétrie hermitienne résulte de l’échange de \(f,g\) et de la conjugaison. Si l’on agrandit \(R\), les termes premier–puissance ajoutés ont des corrélations nulles sur \(\mathcal D_R\), ce qui prouve la compatibilité. \hfill\(\square\)

On peut ainsi écrire une unique forme \(\mathscr W\) sur \(\mathcal D=C_c^\infty(\mathbb R)=\bigcup_R\mathcal D_R\). La compatibilité correspond à la différence complète des contributions ; elle n'exige pas que chaque colonne auxiliaire conserve séparément sa norme lors de l'élargissement du support.

\subsection{Colonnes de transport et identité positive de frontière}
\label{rz:columnas}
La représentation du défaut uniforme du chapitre précédent donne, pour un cycle de longueur \(N\),

\[
\delta_N(T_a)^*\delta_N(T_b)=(I-T_a)^*(I-T_b).
\]
Cette égalité permet de réaliser les différences de translation sur un feuillet de transport normalisé. Pour exposer l'argument sans répéter cette réalisation, on écrira simplement \(\delta_af=(I-T_a)f\). Les évaluations polaires se décomposent en

\[
b_\pm(f)=\frac{\widehat f(i/2)\pm\widehat f(-i/2)}{\sqrt2}.
\]
Définissons deux applications sur \(\mathcal D_R\) :

\[
J_{+,R}f=
\left(
(\sqrt{w_{p,k}}\,\delta_{\ell_{p,k}}f)_{\ell_{p,k}\leq R},
(\sqrt{\mu(a)}\,\delta_af)_{a>0},
b_+(f)
\right),
\]
\[
J_{-,R}f=
\left(
(\sqrt{2w_{p,k}}\,f)_{\ell_{p,k}\leq R},
\sqrt{-c_\Gamma}\,f,
b_-(f)
\right).
\]
Les espaces d'arrivée sont des sommes orthogonales de copies de \(L^2(\mathbb R)\), une intégrale directe dans la composante gamma et une composante scalaire. Lorsque leurs images sont utilisées comme espaces effectifs, on prend les adhérences de \(\operatorname{ran}J_{+,R}\) et de \(\operatorname{ran}J_{-,R}\). L'adhérence d'une image diagonale ne doit pas être confondue avec la somme de toutes les coordonnées de l'espace ambiant.

\HMTresultado{Proposition}{Décomposition exacte de la forme}{rz-num-03-4}
On a

\[
\mathscr W_R(f,g)
=\langle J_{+,R}f,J_{+,R}g\rangle
-\langle J_{-,R}f,J_{-,R}g\rangle.
\]
\textbf{Démonstration.} Pour chaque premier–puissance, la première colonne apporte

\[
w_{p,k}\bigl(2C_{f,g}(0)
-C_{f,g}(\ell_{p,k})-C_{f,g}(-\ell_{p,k})\bigr);
\]
la seconde retire \(2w_{p,k}C_{f,g}(0)\). La composante gamma est précisément l'intégrale de la proposition~\ref{rz-num-03-3}. La composante constante négative produit \(c_\Gamma C_{f,g}(0)\). Enfin,

\[
b_+(f)\overline{b_+(g)}-b_-(f)\overline{b_-(g)}
=h_{f,g}(i/2)+h_{f,g}(-i/2).
\]
La somme de ces identités donne l’expression annoncée. \hfill\(\square\)

Cette décomposition n’établit pas à elle seule la positivité : c’est une différence de deux formes positives. L’identité de frontière du propriétaire démonstratif utilise un espace \(\mathcal K_R\), une projection orthogonale \(P_R\) et un encodeur \(\Xi_R:\mathcal D_R\to\mathcal K_R\) dont les moments doivent satisfaire

\[
\langle\Xi_Rf,\Xi_Rg\rangle
=\langle J_{+,R}f,J_{+,R}g\rangle,
\]
\[
\langle P_R\Xi_Rf,P_R\Xi_Rg\rangle
=\langle J_{-,R}f,J_{-,R}g\rangle.
\]
\HMTresultado{Théorème}{Conséquence des deux moments de frontière}{rz-num-03-5}
Si les deux identités précédentes sont satisfaites sur \(\mathcal D_R\), l'application \(E_R=(I-P_R)\Xi_R\) satisfait

\[
\mathscr W_R(f,g)=\langle E_Rf,E_Rg\rangle.
\]
En particulier, \(\mathscr W_R(f,f)\geq0\).

\textbf{Démonstration.} La décomposition orthogonale de \(\Xi_Rf\) au moyen de \(P_R\) et de \(I-P_R\) donne

\[
\langle\Xi_Rf,\Xi_Rg\rangle
=\langle P_R\Xi_Rf,P_R\Xi_Rg\rangle
+\langle E_Rf,E_Rg\rangle.
\]
On substitue les deux moments et on applique la proposition~\ref{rz-num-03-4}. \hfill\(\square\)

L'importance du théorème réside dans l'identification de la forme arithmétique à une norme de frontière. L'étape qui détermine son application est l'action du codeur et la preuve des deux moments ; l'orthogonalité d'une projection ne les impose pas pour une forme arbitraire. L'annexe de provenance de cette édition enregistre le propriétaire conservé et le point de son intégration matérielle. Cette action n'est pas remplacée par un opérateur défini rétrospectivement à partir de la positivité que l'on veut établir.

Il existe une seconde conservation, utile une fois \(E_R\) construit. La matrice d'incidence

\[
B=\begin{pmatrix}7&2\\2&7\end{pmatrix}=9P_++5P_-
\]
a pour projecteurs orthogonaux

\[
P_\pm=\frac12
\begin{pmatrix}1&\pm1\\\pm1&1\end{pmatrix}.
\]
Sa normalisation donne \(A=P_++qP_-\), avec \(q=5/9\). Si \(D=\sqrt{1-q^2}\,P_-\), alors \(A^*A+D^*D=I\). Par itération,

\[
\|x\|^2=\sum_{n=0}^{M-1}\|DA^nx\|^2+\|A^Mx\|^2.
\]
Sur le feuillet \(P_-x=x\), le terme terminal est \(q^{2M}\|x\|^2\) ; par conséquent, l'application

\[
Lx=(DA^nx)_{n\geq0}
\]
est isométrique sur ce feuillet. Le télescopage explique la conservation au moyen de la distribution des pertes ; l’appartenance de \(E_Rf\) au feuillet doit provenir de la construction précédente. Éteindre un terme terminal ne sélectionne pas rétrospectivement la forme dont on était parti.

\subsubsection*{Codeur interne des histoires et ses deux moments finis}
Le transport corrélatif du TPK permet déjà de construire un codeur avant de définir son résidu. Fixons un horizon \(N\geq1\). Pour chaque niveau, soit \(\mathscr H_k\) la somme orthogonale des deux fibres de feuillet sur les histoires complètes de ce niveau. Si \(d(x)\) compte les émissions sortantes d'une histoire \(x\), le transport agit par

\[
U_k\iota_xv=
\sum_{\alpha:x\to y}\frac1{\sqrt{d(x)}}\,\iota_yU_\alpha v.
\]
On considère ici la réalisation isométrique du transport de feuillet : \(U_\alpha^*U_\alpha=I\). Des histoires distinctes ont des descendants distincts car la route complète est conservée. Par conséquent,

\[
\left.U_k^*U_k\right|_{\mathscr H_x}
=\frac1{d(x)}\sum_{\alpha:x\to y}U_\alpha^*U_\alpha=I.
\]
La cohérence des poids et la conservation de la mémoire sont développées dans la section~\ref{lectura:manuscrito-05} et appliquées au transport spectral dans la sous-section~\ref{rz:refinamiento}. L'isométrie de chaque transport de feuillet est une condition suffisante explicite pour cette égalité ; la normalisation des poids, prise isolément, ne la démontre pas.

Soient \(P_k,Q_k\) les projecteurs orthogonaux des feuillets, avec \(P_k+Q_k=I\). Les parties du transport qui conservent le feuillet et qui en changent sont

\[
A_k=P_{k+1}U_kP_k+Q_{k+1}U_kQ_k,
\qquad
\eta_k=Q_{k+1}U_kP_k+P_{k+1}U_kQ_k.
\]
Le calcul par blocs donne exactement

\[
A_k^*A_k+\eta_k^*\eta_k
=P_kU_k^*U_kP_k+Q_kU_k^*U_kQ_k=I.
\]
La première égalité est valable pour tout transport borné. La seconde utilise la réalisation isométrique indiquée. En général, le premier membre est la partie diagonale par feuillets du Gram de \(U_k\), et pas nécessairement son Gram complet.

On définit \(F_{0,0}=I\), \(F_{0,k}=A_{k-1}\cdots A_0\) pour \(k\geq1\), et \(T_N=F_{0,N}^*F_{0,N}\). L’encodeur interne a pour domaine et pour action

\[
\Xi_N^{\mathrm{int}}:\mathscr H_0\longrightarrow
\mathscr K_N^{\mathrm{int}}
:=\left(\bigoplus_{k=0}^{N-1}\mathscr H_{k+1}\right)\oplus\mathscr H_N,
\]
\[
\Xi_N^{\mathrm{int}}v
=\bigl(\eta_kF_{0,k}v\bigr)_{0\leq k<N}\oplus F_{0,N}v.
\]
Le dernier terme reste séparé même lorsqu’il coïncide comme espace avec celui d’une sortie antérieure. Soit \(\Pi_N^{\mathrm{term}}\) la projection orthogonale sur ce terme terminal. On obtient les deux moments internes

\[
(\Xi_N^{\mathrm{int}})^*\Xi_N^{\mathrm{int}}=I,
\qquad
(\Xi_N^{\mathrm{int}})^*\Pi_N^{\mathrm{term}}\Xi_N^{\mathrm{int}}=T_N.
\]
\textbf{Démonstration.} L'identité par blocs implique

\[
F_{0,k}^*F_{0,k}-F_{0,k+1}^*F_{0,k+1}
=F_{0,k}^*\eta_k^*\eta_kF_{0,k}.
\]
La sommation de \(k=0\) à \(N-1\) produit \(I=\sum_{k<N}F_{0,k}^*\eta_k^*\eta_kF_{0,k}+T_N\), qui est le premier Gram annoncé. La projection terminale ne retient que \(F_{0,N}v\) ; son Gram est donc \(T_N\). Ainsi, le complément \((I-\Pi_N^{\mathrm{term}})\Xi_N^{\mathrm{int}}\) enregistre les sorties de feuillet et a pour Gram \(I-T_N\). Tous ces opérateurs proviennent de la même collection de transports ; le codeur n'est pas défini à partir d'un résidu dont la positivité serait supposée au préalable. \hfill\(\square\)

Ce résultat retrouve un codeur effectif avec ses deux moments, dans le constructeur conjoint qui conserve les quatre autres lectures du continu. Son application aux colonnes complétées de cette section possède un contenu supplémentaire précis : il faudrait identifier une action préalable \(R_R:\mathcal D_R\to\mathscr H_0\), avec la compatibilité des horizons exigée par le domaine des fonctions tests, pour laquelle

\[
\langle R_Rf,R_Rg\rangle
=\langle J_{+,R}f,J_{+,R}g\rangle,
\]
\[
\langle F_{0,N}R_Rf,F_{0,N}R_Rg\rangle
=\langle J_{-,R}f,J_{-,R}g\rangle.
\]
La composition \(\Xi_N^{\mathrm{int}}R_R\) aurait alors les moments du théorème~\ref{rz-num-03-5}. Cette identification de \(R_R\) aux colonnes premier–puissance, gamma et polaire n’est pas réunie ici. La construction interne précédente l’est bien, avec son action et sa preuve ; elle n’est pas présentée comme une démonstration de la positivité globale de \(\mathscr W\). On ne supprime pas non plus \(T_N\), et l’on n’affirme pas que sa limite soit zéro.

\subsubsection*{Compatibilité de l'horizon et type de projection}
La projection terminale \(\Pi_N^{\mathrm{term}}\) sélectionne le dernier terme du codeur d'histoires. L'espérance uniforme \(P_R\) du théorème~\ref{rz-num-03-5} sélectionne, en revanche, la composante uniforme d'un calendrier de neuf positions. L'égalité de leurs fonctions dans une composition doit être démontrée au moyen de l'application concrète entre ces espaces ; toutes deux sont des projections orthogonales, mais cette propriété commune n'identifie pas leurs images.

La dépendance vis-à-vis de l'horizon impose également une identité vérifiable. Soient \(R\) une même application et \(J_-\) une colonne fixe. Si, pour deux horizons consécutifs, on exige

\[
R^*T_NR=J_-^*J_-=R^*T_{N+1}R,
\]
alors

\[
\eta_NF_{0,N}R=0.
\]
\textbf{Démonstration.} En soustrayant les égalités et en utilisant l’identité télescopique, on obtient

\[
0=R^*(T_N-T_{N+1})R
=(\eta_NF_{0,N}R)^*(\eta_NF_{0,N}R).
\]
Son évaluation sur chaque vecteur implique l’annulation indiquée. Si l’égalité des moments est maintenue avec le même \(R\) dans une suite non bornée d’horizons et si \(T_N\) tend fortement vers zéro, on obtient en outre \(J_-=0\). \hfill\(\square\)

Par conséquent, une identification avec une colonne négative non nulle doit spécifier son horizon ou sa collection compatible d'applications \(R_{R,N}\), ou conserver le terme terminal limite qui correspond effectivement. La contraction d'un résidu de sortie ne peut être transférée, sans cette distinction, à la colonne complète représentant le second moment. Cette précision maintient intact le codeur récupéré et détermine la compatibilité supplémentaire que doit vérifier sa réalisation analytique.

\subsection{Représentation des translations}
\label{rz:traslaciones}
Supposons établie sur tout \(\mathcal D\) une identité positive

\[
\mathscr W(f,g)=\langle Ef,Eg\rangle
\]
compatible avec les domaines de support. L'espace spectral est construit à partir de la forme elle-même :

\[
\mathcal N=\{f\in\mathcal D:\mathscr W(f,f)=0\},\qquad
\mathcal H_{\mathscr W}
=\overline{\mathcal D/\mathcal N}.
\]
L’inégalité de Cauchy–Schwarz pour une forme positive prouve que le produit scalaire du quotient est bien défini. L’application \([f]\mapsto Ef\) se prolonge en un isomorphisme unitaire de \(\mathcal H_{\mathscr W}\) sur \(\overline{\operatorname{ran}E}\).

\HMTresultado{Proposition}{Évolution unitaire et générateur}{rz-num-03-6}
Les translations induisent un groupe unitaire fortement continu

\[
U_t[f]=[T_tf].
\]
Son générateur autoadjoint \(H\), avec la convention \(U_t=e^{itH}\), agit sur les classes de fonctions tests par

\[
H[f]=[i f'].
\]
\textbf{Démonstration.} Les corrélations sont invariantes sous la translation commune de leurs deux arguments. De plus, \(\widehat{T_tf}(z)=e^{-izt}\widehat f(z)\) ; dans \(h_{T_tf,T_tg}(z)\), les facteurs se compensent également pour \(z\) complexe. Ainsi, \(\mathscr W(T_tf,T_tg)=\mathscr W(f,g)\), et la translation préserve \(\mathcal N\). Elle définit une isométrie sur le quotient, inversée par \(U_{-t}\), qui se prolonge à l'espace complet.

Pour \(t\to0\), \(T_tf-f\) tend vers zéro dans toutes les semi-normes des fonctions tests sur un compact commun. La continuité prouvée pour \(\mathscr W\) implique la continuité forte sur le quotient dense ; l’unitarité la prolonge à tout l’espace. Le théorème de Stone fournit un générateur autoadjoint. La dérivée en \(t=0\) est \([-f']\), qui coïncide avec \(iH[f]\). \hfill\(\square\)

Les fonctions tests forment en outre un cœur du générateur. On peut le voir en régularisant un vecteur au moyen de \(\int\chi(t)U_t\,dt\), avec \(\chi\in C_c^\infty(\mathbb R)\). Cette régularisation et sa dérivée sont des opérateurs bornés, contrôlés par \(\|\chi\|_1\) et \(\|\chi'\|_1\). On approche d'abord les vecteurs régularisés par des classes de fonctions tests, puis on approche l'identité par une collection de régularisateurs. L'approximation obtenue est dans la norme du graphe \(\|x\|+\|Hx\|\), ce qui prouve l'affirmation.

La représentation a été construite à partir d'une forme positive et de ses translations. Cet argument n'utilise comme entrées ni une liste d'ordonnées de zéros ni une attribution de zéros aux tours du calendrier.

\subsection{Identification analytique et multiplicités}
\label{rz:identificacion}
La forme explicite précédente se reconnaît comme la forme complétée de Weil. Avec

\[
\Phi_f(s)=\widehat f\bigl(i(s-1/2)\bigr),
\qquad \rho^\sharp=1-\overline\rho,
\]
la formule explicite fournit son expression spectrale

\[
\mathscr W(f,g)=\sum_\rho
\Phi_f(\rho)\overline{\Phi_g(\rho^\sharp)},
\]
avec les zéros non triviaux de la fonction complétée et leurs multiplicités. La formulation globale du critère de Weil affirme que la positivité pour toutes les fonctions tests équivaut à \(\Re\rho=1/2\) pour tous ces zéros. On utilise le critère dans la formulation de Bombieri~\cite{bombieri2000}. La condition globale ne peut pas être remplacée par la positivité d'un nombre fini de matrices.

\subsubsection*{Dictionnaire des fonctions tests Mellin--Fourier}
\label{rz:mellin-fourier}
Pour expliciter la normalisation de cette reconnaissance, soit
\(\mathcal D=C_c^\infty(\mathbb R;\mathbb C)\). L’application
\[
 (\mathcal T_{\log}f)(x)=x^{-1/2}f(\log x),\qquad x>0,
\]
est une bijection sur \(C_c^\infty((0,\infty);\mathbb C)\), d'inverse
\[
 (\mathcal T_{\log}^{-1}u)(t)=e^{t/2}u(e^t).
\]
En effet, le changement exponentiel envoie les supports compacts sur des compacts
séparés de zéro et de l'infini. Les formules sont inverses et leurs facteurs
sont lisses. En définissant
\(\widehat u_{\rm Mel}(s)=\int_0^\infty u(x)x^{s-1}\,dx\),
la substitution \(x=e^t\) donne, pour tout \(s\in\mathbb C\),
\[
 \widehat{\mathcal T_{\log}f}_{\rm Mel}(s)
 =\int_{\mathbb R}f(t)e^{(s-1/2)t}\,dt
 =\widehat f\bigl(i(s-1/2)\bigr)=\Phi_f(s).
\]
La conjugaison est conservée exactement :
\[
 \widehat{\overline{\mathcal T_{\log}g}}_{\rm Mel}(1-s)
 =\overline{\Phi_g(1-\overline s)}.
\]
Ainsi, les deux facteurs du terme spectral coïncident avec ceux de la
formule explicite en coordonnées multiplicatives, sans facteur résiduel
et en comptant les mêmes multiplicités. La somme converge absolument :
l’intégration par parties donne une décroissance rapide des transformées
dans les bandes verticales et le dénombrement des zéros est \(O(T\log T)\).
La bijection transporte le quantificateur sur toutes les fonctions tests, non sur une
sélection d’entre elles. La formule explicite et l’implication difficile du
critère sont les antécédents cités ; le changement de variables précédent
démontre leur correspondance de domaines et de normalisations, non une nouvelle
preuve de positivité. Ce domaine compact n’est pas déclaré stable sous
le Cayley intégral : sa réalisation ultérieure utilise l’espace exponentiel
de la sous-section~\ref{tm-add-m4}.

Une fois ces antécédents satisfaits, en écrivant \(\rho=1/2+i\gamma\), il vient

\[
\mathscr W(f,g)
=\sum_{\gamma\in\Gamma_\xi}
m_\gamma\widehat f(-\gamma)\overline{\widehat g(-\gamma)}.
\]
Ici, \(\Gamma_\xi\) contient les ordonnées distinctes et \(m_\gamma\) est la multiplicité du zéro correspondant. Soit

\[
\mu_\xi=\sum_{\gamma\in\Gamma_\xi}m_\gamma\delta_\gamma.
\]
L'application

\[
\mathcal F_\xi[f](\gamma)=\widehat f(-\gamma)
\]
est isométrique de \(\mathcal H_{\mathscr W}\) dans \(L^2(\Gamma_\xi,\mu_\xi)\). Dans celle-ci, \(U_t\) agit comme multiplication par \(e^{it\gamma}\) et \(H\) comme multiplication par \(\gamma\).

\HMTresultado{Proposition}{Exhaustivité de la représentation atomique}{rz-num-03-7}
L'application \(\mathcal F_\xi\) est surjective.

\textbf{Démonstration.} Son image fermée est invariante sous les multiplications par \(e^{it\gamma}\) pour \(t\) positif et négatif ; elle est donc réductrice. Pour une ordonnée fixée \(\gamma_0\), les moyennes

\[
\frac1{2T}\int_{-T}^T e^{-it\gamma_0}U_t\,dt
\]
convergent fortement vers le projecteur sur l'atome \(\gamma_0\). Cela se vérifie d'abord coordonnée par coordonnée, puis par convergence dominée dans l'espace des suites. Il existe \(f\in C_c^\infty\) avec \(\widehat f(-\gamma_0)\neq0\) ; par exemple, il suffit de moduler une fonction d'intégrale non nulle. L'image contient ainsi le vecteur de cet atome. Comme elle contient tous les vecteurs à support fini et est fermée, elle coïncide avec l'espace complet. \hfill\(\square\)

La mesure conserve la multiplicité arithmétique comme poids. Il convient de préciser sa relation avec la multiplicité d'un opérateur : l'espace scalaire \(L^2(\Gamma_\xi,\mu_\xi)\) possède un sous-espace propre de dimension un pour chaque atome, bien que \(m_\gamma>1\). Ainsi, si une fonction \(F\) produit un opérateur à trace,

\[
\operatorname{Tr}(F(H))=\sum_{\gamma\in\Gamma_\xi}F(\gamma),
\]
tandis que la somme qui conserve les poids des zéros est

\[
\int F(\gamma)\,d\mu_\xi(\gamma)
=\sum_{\gamma\in\Gamma_\xi}m_\gamma F(\gamma).
\]
Cette seconde somme peut se représenter comme une trace pondérée \(\operatorname{Tr}(M_mF(H))\), lorsque le produit est à trace, avec \(M_m\) l'opérateur de multiplication par \(m_\gamma\). Une autre construction, distincte, serait une fibre de dimension \(m_\gamma\) en chaque atome. Ces deux représentations ne sont pas échangées sans le déclarer.

\subsection{Les deux cercles et l’orientation spectrale}
\label{rz:circulos}
Le premier cercle est une coordonnée compacte de la droite spectrale. Pour un nombre réel \(\gamma\), définissons

\[
\tau_\gamma=\frac{\gamma+i/2}{\gamma-i/2}.
\]
Son module est un et \(\tau_\gamma\neq1\). Les formules inverses sont

\[
\gamma=\frac{i}{2}\frac{\tau_\gamma+1}{\tau_\gamma-1},
\qquad
\gamma=\frac12\cot\frac{\theta}{2}
\quad\text{si}\quad\tau_\gamma=e^{i\theta}.
\]
Si \(s=1/2+i\gamma\), on a également

\[
\tau_\gamma=\frac{s-1}{s},\qquad s=\frac1{1-\tau_\gamma}.
\]
L’ordonnée n’a pas été perdue lors du passage à la phase, pourvu que le point du cercle et la convention d’orientation soient maintenus.

Pour le générateur autoadjoint précédent, le calcul fonctionnel définit l'opérateur unitaire

\[
C=(H+iI/2)(H-iI/2)^{-1}.
\]
L'égalité \(C^*C=I\) se déduit du fait que les deux facteurs sont des fonctions du même opérateur autoadjoint et que \(|(\gamma+i/2)/(\gamma-i/2)|=1\). Le point \(1\) ne correspond pas à une ordonnée finie, mais peut appartenir au spectre de \(C\) comme point d'accumulation lorsque \(H\) est non borné.

Le second cercle est le dual de la mémoire entière : à une phase \(\tau\) est associé le caractère

\[
\chi_\tau(n)=\tau^n,\qquad n\in\mathbb Z.
\]
La connexion nonadique ramène la phase visible à sa position initiale et augmente d’une unité la mémoire d’un tour. La réalisation spectrale assigne à cette unité l’action de \(C\) ; sur un vecteur propre d’ordonnée \(\gamma\), la mémoire de \(n\) tours agit par \(\tau_\gamma^n\). Le premier cercle encode l’ordonnée ; le second exprime comment sa phase agit sur une mémoire additive. Ce sont deux fonctions mathématiques liées, non deux listes devant être identifiées par leur position ordinale.

\subsection{Raffinement cylindrique avec mémoire}
\label{rz:refinamiento}
La connexion doit agir sur la structure qui est raffinée. Dans l'arbre commun du TPK, un sommet est une histoire finie d'émissions admissibles ; deux histoires qui atteignent le même état terminal restent distinctes lorsque leurs registres de transport diffèrent. Si \(\operatorname{Out}(x)\) est l'ensemble fini des émissions sortantes, comptées avec leur multiplicité, on écrit

\[
d(x)=|\operatorname{Out}(x)|\geq1.
\]
Le prolongement neutre conserve la non-vacuité des niveaux. Le poids local déclaré par l'arbre est \(1/d(x)\). Pour une histoire \(h=(\epsilon_1,\ldots,\epsilon_k)\), d'états successifs \(x_0,\ldots,x_k\), son poids est

\[
\mu_k(h)=\prod_{j=0}^{k-1}\frac1{d(x_j)}.
\]
La somme des poids de tous ses successeurs est

\[
\sum_{\epsilon\in\operatorname{Out}(x_k)}
\mu_{k+1}(h\epsilon)
=\mu_k(h)\sum_{\epsilon\in\operatorname{Out}(x_k)}\frac1{d(x_k)}
=\mu_k(h).
\]
En partant de la racine de poids un, cette identité démontre par récurrence la normalisation et la cohérence projective. Si une symétrie réversible transporte bijectivement les émissions sortantes avec leurs multiplicités, elle conserve \(d(x)\), tous ces poids et leurs mesures de cylindres. Ainsi, la mesure utilisée ci-après provient de la règle de ramification déclarée, avant d'attribuer un opérateur à ses fibres. Elle n'est pas confondue avec une distribution uniforme sur des états terminaux qui identifierait des histoires différentes.

De manière générale, soit \(\mathcal Z_k\) un niveau fini d'états cylindriques et soit \(\pi_k:\mathcal Z_{k+1}\to\mathcal Z_k\) son application de troncature. On suppose donnée une collection de poids positifs compatibles,

\[
\mu_k(z)=\sum_{\pi_k(z')=z}\mu_{k+1}(z'),
\qquad \sum_{z\in\mathcal Z_k}\mu_k(z)=1.
\]
Les poids appartiennent à la mesure d’états qui est transportée. La forme d’un cylindre numérique visible ne suffit pas à elle seule à les reconstruire.

Soient \(q_k:\mathcal Z_k\to\mathbb Z\) les coordonnées de mémoire. L'incrément d'un prolongement est

\[
m(z',z)=q_{k+1}(z')-q_k(z).
\]
En concaténant des prolongements, on vérifie exactement la loi de cocycle :

\[
m(z'',z)=m(z'',z')+m(z',z).
\]
On définit

\[
w(z',z)=
\sqrt{\frac{\mu_{k+1}(z')}{\mu_k(z)}},
\qquad
\mathcal K_k=\ell^2(\mathcal Z_k)\otimes\mathcal H_{\mathscr W}.
\]
Le transport avec mémoire est

\[
\widetilde U_k(e_z\otimes v)
=\sum_{\pi_k(z')=z}w(z',z)e_{z'}\otimes C^{m(z',z)}v.
\]
\HMTresultado{Proposition}{Isométrie et composition du transport}{rz-num-03-8}
Chaque \(\widetilde U_k\) est isométrique. La composition sur plusieurs niveaux a la même forme, avec le quotient des poids entre les niveaux extrêmes et l’accroissement total de mémoire.

\textbf{Démonstration.} Deux prédécesseurs distincts ont des ensembles disjoints de descendants. Pour un prédécesseur fixé, les puissances de \(C\) sont unitaires et \(\sum_{z'\mapsto z}w(z',z)^2=1\). Ce sont exactement les identités qui conservent les produits scalaires des vecteurs de base. Dans une composition, le produit des poids est télescopique :

\[
\sqrt{\frac{\mu_{k+2}(z'')}{\mu_{k+1}(z')}}
\sqrt{\frac{\mu_{k+1}(z')}{\mu_k(z)}}
=\sqrt{\frac{\mu_{k+2}(z'')}{\mu_k(z)}}.
\]
Les puissances de \(C\) additionnent leurs exposants par la loi de cocycle. La répétition prouve l'affirmation pour tout nombre fini de niveaux. \hfill\(\square\)

La même identité admet une lecture par changement de représentation. Si \(U_k\) désigne le transport sans mémoire et

\[
G_k(e_z\otimes v)=e_z\otimes C^{q_k(z)}v,
\]
alors

\[
\widetilde U_k=G_{k+1}(U_k\otimes I)G_k^*.
\]
La coordonnée de mémoire n'est pas perdue lors de la formulation de l'isométrie : elle apparaît aux deux extrémités du transport. Cette égalité montre également que la représentation tordue utilise un opérateur \(C\) déjà spécifié. À elle seule, elle n'identifie pas tout transport natif antérieur à cet opérateur ; l'identification exige de composer leurs applications et de vérifier leurs actions sur le même domaine.

\HMTresultado{Proposition}{Retour nonadique sur la section compatible}{rz-num-03-9}
On définit

\[
\Omega_k=\sum_z\sqrt{\mu_k(z)}e_z,\qquad
J_kv=\Omega_k\otimes v.
\]
Si toute trajectoire du retour de \(k\) à \(k+9\) a un incrément de mémoire égal à un, alors

\[
\widetilde U_{\Gamma,k}J_k=J_{k+9}C.
\]
Pour \(r\) tours consécutifs,

\[
\widetilde U_{\Gamma,k+9(r-1)}\cdots
\widetilde U_{\Gamma,k}J_k=J_{k+9r}C^r.
\]
\textbf{Démonstration.} En multipliant le poids de chaque descendant par le coefficient \(\sqrt{\mu_k(z)}\) de son prédécesseur, on obtient \(\sqrt{\mu_{k+9}(z')}\). La condition sur l'incrément permet de mettre le même opérateur \(C\) en facteur de chaque terme. On obtient la première égalité. La seconde en découle par récurrence. \hfill\(\square\)

Les indices de niveau sont essentiels. Le retour est celui de la phase et la section se prolonge à un autre niveau ; la formule ne réinitialise pas l'état dans l'espace précédent.

La condition de mémoire de cette proposition possède une action concrète dans la chronologie du TPK. Pour une phase \(r\in\mathbb Z/9\mathbb Z\), de représentant \(\bar r\in\{0,\ldots,8\}\), et une profondeur de mémoire \(w\in\mathbb Z_{\geq0}\), chaque émission met à jour

\[
\sigma_9(w,r)=
\left(w+\left\lfloor\frac{\bar r+1}{9}\right\rfloor,r+[1]_9\right).
\]
Après \(n\) émissions, le nombre de franchissements du représentant \(8\) au représentant \(0\) est \(\lfloor(\bar r+n)/9\rfloor\). Par conséquent,

\[
\sigma_9^n(w,r)=
\left(w+\left\lfloor\frac{\bar r+n}{9}\right\rfloor,r+[n]_9\right),
\qquad
\sigma_9^9(w,r)=(w+1,r).
\]
Prendre \(q_k=w_k\) retrouve le cocycle précédent et l'incrément un de chaque tour. Même une émission neutre dans les feuillets d'APP enregistre une extension dans cette chronologie. Le calcul n'impose pas par lui-même l'opérateur spectral \(C\) : il détermine la mémoire sur laquelle agit la réalisation, une fois son entrelacement établi.

\subsection{Calendrier nonadique et contraction d’un tour}
\label{rz:calendario}
Le calendrier fini est une représentation de ce transport. Sur \(\mathbb C^9\otimes\mathcal H_{\mathscr W}\), soit

\[
S_C=\sum_{r=1}^{8}|r+1\rangle\langle r|\otimes I
+|1\rangle\langle9|\otimes C.
\]
L'unitarité de \(C\) rend \(S_C\) unitaire, et le tour complet satisfait

\[
S_C^9=I_{\mathbb C^9}\otimes C.
\]
Chaque vecteur traverse exactement une fois la jonction orientée qui contient \(C\) ; c’est pourquoi \(C^9\) n’apparaît pas.

Soient \(u=9^{-1/2}(1,\ldots,1)\) et \(\iota v=u\otimes v\). La compression uniforme d'un pas est

\[
A_C=\iota^*S_C\iota=\frac{8I+C}{9}.
\]
Un calcul direct donne

\[
I-A_C^*A_C
=\frac8{81}(I-C)^*(I-C).
\]
La différence par rapport à un transport unitaire comprimé mesure la composante qui quitte l'état uniforme. Ce n'est pas une perte d'unitarité du calendrier complet.

La contraction \(q=5/9\) du feuillet complémentaire et le calendrier sont des opérations distinctes. Si \(q\) correspond à un tour nonadique, l’opérateur par tour est

\[
M_{\mathrm{vuelta}}=qS_C^9
=q(I\otimes C).
\]
Elle peut être répartie uniformément entre les neuf pas au moyen de

\[
M_{\mathrm{paso}}=q^{1/9}S_C,\qquad
M_{\mathrm{paso}}^9=M_{\mathrm{vuelta}}.
\]
En revanche, prendre \(qS_C\) comme opérateur de chaque pas produit \(q^9(I\otimes C)\) lorsque le tour est achevé. Les deux normalisations décrivent des contractions différentes et ne doivent pas être identifiées. Sur une ordonnée \(\gamma\), \(r\) tours du premier régime produisent

\[
q^r\tau_\gamma^r.
\]
Le module enregistre l'extinction du terme terminal ; la phase enregistre l'action sur la mémoire. Connaître le module seul ne détermine pas les ordonnées spectrales.

\subsection{Matrices de moments et approximation spectrale}
\label{rz:galerkin}
La réalisation précédente permet de construire des approximations qui partent de fonctions tests, non de valeurs spectrales assignées. Choisissons \(f_1,\ldots,f_d\in\mathcal D\) dont les classes sont linéairement indépendantes dans \(\mathcal H_{\mathscr W}\). Avec \(e_j=[f_j]\), on définit

\[
G_{ij}=\langle e_j,e_i\rangle,\qquad
K_{ij}=\langle He_j,e_i\rangle,\qquad
Q_{ij}=\langle He_j,He_i\rangle.
\]
La convention d'indices fait que, pour \(x=\sum_jv_je_j\), on a \(\|x\|^2=v^*Gv\). \(G\) est hermitienne définie positive et \(K,Q\) sont hermitiennes. Si les classes initiales étaient dépendantes, il faut retirer le noyau de \(G\) avant de l'inverser.

\HMTresultado{Proposition}{Approximation de Galerkin et résidu}{rz-num-03-10}
Le problème généralisé

\[
Kv=\gamma Gv
\]
équivaut au problème hermitien ordinaire de \(\widehat H=G^{-1/2}KG^{-1/2}\). Pour tous \(\gamma\in\mathbb R\) et \(v\neq0\), le résidu spectral du vecteur correspondant satisfait

\[
\frac{\|(H-\gamma)x\|^2}{\|x\|^2}
=\frac{v^*(Q-2\gamma K+\gamma^2G)v}{v^*Gv}.
\]
En particulier,

\[
\operatorname{dist}(\gamma,\sigma(H))
\leq\frac{\|(H-\gamma)x\|}{\|x\|}.
\]
\textbf{Démonstration.} Le changement \(w=G^{1/2}v\) transforme l’équation généralisée en \(\widehat Hw=\gamma w\). Le développement de la norme du résidu donne l’identité quadratique car \(H\) est symétrique sur les vecteurs utilisés. La dernière inégalité se déduit de la mesure spectrale de \(H\) :

\[
\|(H-\gamma)x\|^2
=\int |\lambda-\gamma|^2\,d\langle E_H(\lambda)x,x\rangle.
\]
L'intégrande est minoré par le carré de la distance au spectre. \hfill\(\square\)

La matrice de phase

\[
\widehat C=(\widehat H+iI/2)(\widehat H-iI/2)^{-1}
\]
est unitaire pour le produit euclidien. Si l'on utilise plutôt \(G^{-1}K\) dans la base d'origine, l'unitarité correspondante s'exprime par rapport à \(G\). Ce sont deux bases du même problème, et non deux normes interchangeables sans transformation.

Un petit résidu fournit une approximation du spectre de l'opérateur déjà construit. À lui seul, il ne compte pas toutes ses valeurs propres, ne démontre pas la positivité globale de \(\mathscr W\) et ne certifie pas les décimales affectées par les erreurs de quadrature. Pour une borne validée, les entrées de \(G,K,Q\) doivent être accompagnées de bornes d'erreur et d'une borne inférieure positive de la forme \(v^*Gv\).

\subsection{Résolvante compacte et contrôles d’exhaustivité}
\label{rz:resolvente}
Dans la réalisation atomique, les ordonnées forment un ensemble discret et tendent vers l’infini. Par conséquent,

\[
R=(H-iI)^{-1}
\]
est compact : dans la base atomique, il agit par \((\gamma-i)^{-1}\), qui tend vers zéro. On le retrouve également directement à partir du groupe des translations au moyen de

\[
R=i\int_0^\infty e^{-t}U_{-t}\,dt.
\]
L’intégrale sur un intervalle fini est définie au sens fort. Sa queue à partir de \(T\) a une norme opératorielle au plus égale à \(e^{-T}\), de sorte que les intégrales tronquées convergent en norme. Le calcul fonctionnel ramène l’identité à \(i/(1+i\gamma)=(\gamma-i)^{-1}\).

\HMTresultado{Proposition}{Convergence en norme des compressions}{rz-num-03-11}
Si \(P_d\) sont des projecteurs orthogonaux de rang fini avec \(P_d\to I\) fortement, alors

\[
\|R-P_dRP_d\|\longrightarrow0.
\]
\textbf{Démonstration.} La convergence forte de \(P_d\), jointe à la borne \(\|P_d\|\leq1\), est uniforme sur tout compact. On l’applique à l’adhérence compacte de l’image de la boule unité par \(R\), ce qui donne \(\|(I-P_d)R\|\to0\). En appliquant le même argument à \(R^*\), on obtient \(\|R(I-P_d)\|\to0\). La décomposition

\[
R-P_dRP_d=(I-P_d)R+P_dR(I-P_d)
\]
conclut la preuve. \hfill\(\square\)

Cette convergence permet des contrôles par des courbes de résolvante. Si une courbe fermée \(\mathcal C\) sépare un groupe de valeurs propres non nulles et si la norme de l’erreur de compression est inférieure à l’inverse d’une borne de la résolvante approchée sur \(\mathcal C\), la série de Neumann conserve l’inversibilité le long de la courbe. Les projecteurs de Riesz des deux opérateurs ont alors le même rang. C’est un contrôle d’exhaustivité dans une région spectrale délimitée, supplémentaire au résidu d’un vecteur.

La multiplicité comptée par ce rang est celle de l'opérateur représenté. Dans l'espace scalaire de la sous-section~\ref{rz:identificacion}, elle ne doit pas être remplacée par le poids \(m_\gamma\) sans introduire expressément l'observable pondéré.

\subsection{Des identités internes à l’assemblage analytique}
\label{rz:alcance}
Sont réunies ici les preuves de la normalisation de la chaleur et de Mellin, la construction explicite de la forme complétée, sa décomposition par colonnes, le transport cylindrique avec mémoire, les deux cercles, le calendrier et les contrôles spectraux. La construction arithmétique des irréductibles précède leurs apparitions dans les poids premier–puissance. La mémoire provient des prolongements compatibles ; elle n'est pas remplacée par un chiffre de phase.

La sous-section~\ref{rz:columnas} réunit en outre un encodeur interne fini construit à partir des transports d’histoires, ainsi que la preuve de ses moments \(I\) et \(T_N\) dans la réalisation isométrique. L’application du théorème~\ref{rz-num-03-5} à l’encodeur natif de la forme complétée conserve une identification plus spécifique : l’action préalable \(R_R\) sur les fonctions tests et l’égalité de ces moments avec ceux de \(J_{+,R}\) et \(J_{-,R}\). Le propriétaire de la forme complétée formule ces identités et développe leurs conséquences. Cette rédaction conserve ces conséquences avec leurs antécédents explicites et la composition interne effectivement retrouvée ; elle ne remplace pas le lien restant par une formulation nominale. La distinction concerne la démonstration réunie dans cet article, non une déclaration d’absence du résultat dans l’ensemble du corpus.

Le lien complet entre toutes les histoires APP–TRIT–TPK et la normalisation par mots du premier chapitre conserve également son obligation d’intégration. Les cinq constructions consubstantielles du continu appartiennent toujours au noyau commun ; ce chapitre expose leur lecture spectrale et ne les redéfinit pas comme des structures indépendantes sélectionnées par des problèmes conventionnels.

L'identification de la forme globale à la formule explicite et au critère de Weil fixe exactement la portée d'une éventuelle conclusion terminale. Les vérifications finies du producteur natif et les approximations spectrales sont des résultats vérifiables distincts. C'est pourquoi cette édition de travail ne présente pas le chapitre comme une démonstration autonome déjà achevée de l'hypothèse de Riemann.

\subsection{Récupération bidirectionnelle et assemblage des lecteurs}
\label{rz:gamma}
L’information gamma récupère la fonction test qui l’a produite. Cette propriété fournit une action explicite entre les colonnes complètes, avant de connaître le signe de leur différence. Le développement suivant étend l’état d’intégration décrit dans les sous-sections~\ref{rz:columnas} et~\ref{rz:alcance} : le connecteur borné global est récupéré ; sa contractivité est étudiée sur la même action et la même image, non sur un autre opérateur.

La construction se situe après les poids premier–puissance, la trace complétée et les transports déjà définis. Elle ne change pas la sélection des irréductibles et ne reçoit pas de zéros en entrée. Les propriétaires de la queue gamma, du domaine basal et du pliage préalable sont conservés intégralement dans le paquet de provenance de cette intégration.

\HMTresultado{Proposition}{Inverse à gauche de la queue et connecteur global}{rz-num-03-12}
Soit \(H_R=L^2(I_R)\), identifié à son prolongement par zéro à la droite, et soit \(M_R\) la restriction orthogonale à \(I_R\). Définissons

\[
a_R=2\int_R^\infty\mu(a)\,da,\qquad
(B_Rf)(a)=\sqrt{\mu(a)}(I-T_a)f\quad(a\geq R).
\]
Alors \(B_R\) est borné et \(B_R^*B_R=a_RI\). Si \(\pi_{\Gamma,\geq R}\) extrait la queue gamma de la colonne positive, l’application

\[
L_R=a_R^{-1}B_R^*\pi_{\Gamma,\geq R}
\]
satisfait \(L_RJ_{+,R}f=f\) pour \(f\in\mathcal D_R\), et \(L_RL_R^*=a_R^{-1}I\) sur \(H_R\). La première identité s’étend au domaine fermé conjoint des lecteurs ; elle n’affirme pas que la colonne gamma complète soit définie sur tout \(L^2(I_R)\). En particulier,

\[
\mathcal C_R=J_{-,R}L_R,\qquad
\mathcal C_RJ_{+,R}=J_{-,R}
\]
est un connecteur borné pour tout \(R>0\).

\textbf{Démonstration.} Pour \(a\geq R\), les supports de \(f,g\in H_R\) et de ses translatées sont séparés, sauf aux extrémités de mesure nulle. Par conséquent

\[
\langle(I-T_a)f,(I-T_a)g\rangle=2\langle f,g\rangle.
\]
L’intégration donne le premier Gram. L’adjoint possède l’action explicite

\[
B_R^*F=M_R\int_R^\infty\sqrt{\mu(a)}(I-T_{-a})F(a)\,da.
\]
L’intégrale converge en norme par Cauchy–Schwarz ; la queue de \(\mu\) est intégrable et \(T_a^*=T_{-a}\). Appliquer l’adjoint à la queue de \(J_{+,R}f\) restitue \(a_Rf\). L’extraction est une coisométrie sur son facteur direct, d’où \(L_RL_R^*=a_R^{-2}B_R^*B_R=a_R^{-1}I\). La colonne négative est bornée sur \(H_R\) : elle comporte un nombre fini de lignes vectorielles et une fonctionnelle polaire continue sur cet intervalle. La composition conclut. \hfill\(\square\)

L'inversion utilise les translations conjuguées avant d'exprimer la fonction récupérée. La source de référence fournit également une version à bande finie : avec \(A_R=[R+1,R+2]\) et \(\nu_R=\int_{A_R}\mu(a)\,da\),

\[
\mathfrak R_RF=\nu_R^{-1}\int_{A_R}
\sqrt{\mu(a)}M_RF(a)\,da,\qquad
\mathfrak R_RD_\Gamma f=f.
\]
L’identité résulte de \(M_RT_af=0\) sur cette bande. Les deux inverses sont antérieurs à une inégalité de Weil. Le premier permet de calculer exactement la norme ambiante de son connecteur.

\HMTresultado{Proposition}{Norme du connecteur et factorisation basale}{rz-num-03-13}
Écrivons

\[
s_R^2=2\sum_{\ell_{p,k}\leq R}w_{p,k}-c_\Gamma,\qquad
\beta_R=2\sinh(R/2)-R.
\]
Alors

\[
\|\mathcal C_R\|^2=\frac{s_R^2+\beta_R}{a_R}.
\]
Par conséquent, si \(0<h<\log2\) et

\[
2M_\Gamma(h)\geq-c_\Gamma+2\sinh(h/2)-h,\qquad
M_\Gamma(h)=\int_h^\infty\mu(a)\,da,
\]
il existe une dilatation qui réalise les deux moments du théorème~\ref{rz-num-03-5} sur toutes les fonctions de \(\mathcal D_h\).

\textbf{Démonstration.} La colonne négative a pour matrice de Gram \(J_-^*J_-=s_R^2I+b_-^*b_-\). La fonctionnelle

\[
b_-(f)=\sqrt2\int_{-R/2}^{R/2}f(x)\sinh(x/2)\,dx
\]
a pour carré de sa norme \(2\int_{-R/2}^{R/2}\sinh^2(x/2)\,dx=\beta_R\). La matrice de Gram a pour norme \(s_R^2+\beta_R\). L'identité précédente \(L_RL_R^*=a_R^{-1}I\) donne \(\mathcal C_R\mathcal C_R^*=a_R^{-1}J_-J_-^*\) et prouve la formule.

Pour \(R=h<\log2\), il n’y a pas de longueurs premier–puissance dans la somme. L’inégalité déclarée prouve \(\|\mathcal C_h\|\leq1\). On peut construire, dans cet ordre,

\[
D_h^{\mathrm{def}}=(I-\mathcal C_h^*\mathcal C_h)^{1/2},
\qquad
\Xi_hf=J_{-,h}f\oplus D_h^{\mathrm{def}}J_{+,h}f.
\]
La somme des Grams de \(\mathcal C_h\) et de son défaut est l’identité. Comme \(\mathcal C_hJ_{+,h}=J_{-,h}\), le premier moment de \(\Xi_h\) est \(J_{+,h}^*J_{+,h}\) ; la projection sur le premier facteur direct a pour moment \(J_{-,h}^*J_{-,h}\). Le théorème~\ref{rz-num-03-5} donne la positivité. La racine du défaut est prise après avoir démontré la contractivité, et non pour remplacer cette preuve.

Si la fenêtre est étendue à \(R\geq h\), chaque canal premier ajouté satisfait \(\|(I-T_\ell)f\|^2=2\|f\|^2\) sur \(H_h\). L'application de cette image dans \(\sqrt2f\) est une isométrie, prolongée par zéro sur le complément de son image fermée. La somme directe de ces lignes avec le connecteur gamma–polaire maintient la contractivité dans le même domaine basal. \hfill\(\square\)

L’extrémité basale est déterminée par

\[
2M_\Gamma(h_*)=-c_\Gamma+2\sinh(h_*/2)-h_*.
\]
Le changement \(t=e^{-a/2}\) donne

\[
M_\Gamma(h)=\operatorname{arctanh}(e^{-h/2})
+\arctan(e^{-h/2}).
\]
Le premier membre de l’équation décroît strictement de l’infini à zéro ; le second est positif et croît strictement pour \(h>0\). Il existe une unique racine positive. Le propriétaire l’évalue à \(h_*=0.08562601762973687068\ldots<\log2\). On peut travailler avec tout \(h<h_*\) certifié par des inégalités, sans avoir besoin des décimales de cette extrémité.

La forme complète conserve également la translation des fonctions tests. Dans les ports polaires, celle-ci agit par

\[
\binom{b_+(T_tf)}{b_-(T_tf)}=
\begin{pmatrix}\cosh(t/2)&\sinh(t/2)\\
\sinh(t/2)&\cosh(t/2)\end{pmatrix}
\binom{b_+(f)}{b_-(f)}.
\]
La matrice conserve la différence de carrés et les autres lignes sont translatées unitairement. Le domaine basal peut donc être situé dans n'importe quel intervalle translaté de longueur \(h\).

\HMTresultado{Proposition}{Unicité sur l’image effective}{rz-num-03-14}
Soit \(\mathcal M_{+,R}=\overline{J_{+,R}\mathcal D_R}\). La restriction du connecteur construit à ce sous-espace est unique parmi les applications bornées possédant cette identité de sortie. De plus,

\[
\|\mathcal C_R|_{\mathcal M_{+,R}}\|\leq1
\quad\Longleftrightarrow\quad
\mathscr W_R(f,f)\geq0\quad(f\in\mathcal D_R).
\]
\textbf{Démonstration.} Deux opérateurs bornés qui coïncident sur l’image dense coïncident sur sa fermeture. L’équivalence résulte de \(\mathcal C_RJ_{+,R}f=J_{-,R}f\) et de l’identité des normes ; l’inégalité se prolonge à la fermeture par continuité. \hfill\(\square\)

La formule globale est donc construite :

\[
\mathscr W_R(f,g)=
\langle J_{+,R}f,(I-\mathcal C_R^*\mathcal C_R)J_{+,R}g\rangle.
\]
Changer la représentation ambiante peut faciliter le calcul de la norme, mais ne permet pas de changer son action sur les fonctions tests. La non-contractivité dans tout l’espace ambiant ne réfute pas la positivité sur l’image complète, où subsistent les corrélations des lignes. La projection sur cette image ne prouve pas non plus automatiquement la borne.

\HMTresultado{Proposition}{Pliage préalable et totalité des produits croisés}{rz-num-03-15}
Une fibre basale étant fixée, soient \(m=9^q\) et \(\ell_j\) les déplacements exprimés par ses routes. Sur \(m\) copies de \(L^2(\mathbb R)\), le pliage et sa représentation uniforme sont

\[
\mathcal F_qx=\sum_{j=1}^mT_{\ell_j}x_j,\qquad
\mathsf S_qx=u_q\otimes\mathcal F_qx=3^qP_q\mathsf T_qx,
\qquad u_q=m^{-1/2}(1,\ldots,1).
\]
Pour tout lecteur linéaire \(J\) défini sur ces fonctions tests,

\[
\|J\mathcal F_qx\|^2
=\sum_{i,j}\langle JT_{\ell_i}x_i,JT_{\ell_j}x_j\rangle.
\]
En outre, \(\mathcal F_q\) est surjective, avec pour inverse à droite

\[
(\mathcal E_qf)_j=m^{-1}T_{-\ell_j}f,\qquad
\mathcal F_q\mathcal E_q=I,\qquad
\|\mathcal E_qf\|^2=m^{-1}\|f\|^2.
\]
\textbf{Démonstration.} Prendre la moyenne des \(m\) routes et multiplier par \(\sqrt m=3^q\) donne l’expression de \(\mathsf S_q\). Développer le produit scalaire conserve tous les croisements. Substituer \(\mathcal E_qf\) produit \(m\) copies de \(f/m\) ; l’unitarité des translations donne la norme. \hfill\(\square\)

Si les successeurs vérifient \(\ell_{jr}=\ell_j+\delta_{jr}\), l'opération \((\mathcal A_qx)_j=\sum_{r=0}^8T_{\delta_{jr}}x_{jr}\) satisfait exactement

\[
\mathcal F_{q+1}=\mathcal F_q\mathcal A_q,\qquad
\mathsf S_{q+1}=\mathcal I_q\mathsf S_q\mathcal A_q,\qquad
\mathcal I_qy=y\otimes u_9.
\]
Les colonnes gamma, première et polaire s'appliquent ainsi à la même somme avant le calcul de leurs normes. Un mélange unitaire postérieur à une somme orthogonale préserve l'énergie diagonale ; il ne fabrique pas de termes croisés omis.

La surjectivité détermine la portée du raffinement. Dans la fibre complète, la positivité de \(\mathscr W(\mathcal F_qx,\mathcal F_qx)\) pour tous les états équivaut à la positivité de \(\mathscr W(f,f)\) pour toutes les fonctions tests. L’implication réciproque utilise \(x=\mathcal E_qf\), dont les composantes restent lisses à support compact. Le pliage préserve le problème et ses croisements : il ne restreint pas les fonctions tests du seul fait d’augmenter les routes. Une restriction à des états atteignables plus spécifiques exige sa construction et sa preuve.

\subsubsection*{Résultat réuni et extension globale}
L’inverse gamma, le lien complet de sortie, le domaine basal, sa translation et le pli antérieur aux lecteurs ont ici une action, un domaine et une démonstration. Ce sont des résultats récupérés du corpus et des formalisations réunies, non des découvertes attribuées rétrospectivement.

Pour la portée globale demandée, il reste à démontrer la borne de cette même sortie sur toute l'image effective, y compris les termes croisés des sommes de fonctions tests locales, de manière compatible avec l'élargissement des fenêtres. L'identité de sortie d'un connecteur et la contractivité d'un autre ne se combinent pas sans démontrer qu'ils représentent la même action. La comparaison avec le codeur interne d'histoires conserve les deux moments spécifiés dans la sous-section~\ref{rz:columnas}.

\subsection{Fenêtres d’horloges irréductibles et raffinement avec termes croisés}
\label{rz:ventanas}
Le produit natif détermine les irréductibles et leurs puissances avant la lecture logarithmique. Les longueurs \(\ell_{p,k}=k\log p\) et les poids \(w_{p,k}=(\log p)p^{-k/2}\) qui apparaissent ci-dessous sont ceux déjà exprimés dans les sections~\ref{lectura:manuscrito-01} et~\ref{lectura:manuscrito-02}. Ils ne sont pas sélectionnés à partir du signe que l'on veut vérifier. Sur cette sortie arithmétique, on peut calculer conjointement une collection de fonctions tests, en maintenant les canaux de la forme complétée et tous leurs produits croisés.

Cette section réunit la densité et le critère de Schur de la source « Lecteur survivant de deux routes APP », et étend le calcul à des intervalles translatés de même largeur. Le choix de neuf fenêtres certifié ci-dessous est une collection analytique explicite ; son cardinal n'identifie pas à lui seul ces fenêtres aux neuf successeurs canoniques d'un état TPK.

\HMTresultado{Proposition}{Noyau complet d'intervalles translatés}{rz-num-03-16}
Pour \(\varepsilon>0\), soient

\[
b_{\varepsilon,t}(x)=\varepsilon^{-1/2}
\mathbf1_{[t-\varepsilon/2,t+\varepsilon/2]}(x),\qquad
\tau_\varepsilon(x)=\left(1-\frac{|x|}{\varepsilon}\right)_+,
\]
\[
a_\varepsilon=\frac{4\sinh(\varepsilon/4)}{\sqrt\varepsilon},
\qquad \lambda_n=2n+\frac12,\qquad
F(x)=\sum_{n\geq0}\frac{e^{-\lambda_nx}}{\lambda_n^2}
\quad(x\geq0).
\]
La forme complète sur ces intervalles dépend de la différence des centres :

\[
\mathscr W(b_{\varepsilon,s},b_{\varepsilon,t})
=K_\varepsilon(s-t),
\]
où

\[
\begin{aligned}
K_\varepsilon(d)={}&2a_\varepsilon^2\cosh(d/2)
+c_\Gamma\tau_\varepsilon(d)\\
&+\frac{2F(|d|)-F(|d-\varepsilon|)-F(|d+\varepsilon|)}{\varepsilon}\\
&-\sum_{p,k}w_{p,k}\bigl[
\tau_\varepsilon(d-\ell_{p,k})+
\tau_\varepsilon(d+\ell_{p,k})\bigr].
\end{aligned}
\]
La somme premier–puissance est finie ; seules contribuent les longueurs \(\ell_{p,k}<|d|+\varepsilon\).

\textbf{Démonstration.} Le recouvrement de deux intervalles normalisés est la corrélation triangulaire \(\tau_\varepsilon\). L'insertion des translations de chaque horloge fournit les deux termes premiers. Les intégrales polaires de \(b_{\varepsilon,t}\) sont \(a_\varepsilon e^{t/2}\) et \(a_\varepsilon e^{-t/2}\). Leur produit croisé donne \(a_\varepsilon^2(e^{(s-t)/2}+e^{(t-s)/2})\), le terme polaire de la formule. En particulier, la dépendance à la différence des centres appartient à la forme complète, bien que chaque port polaire pris séparément change sous l'effet d'une translation.

Le développement \(\mu(a)=\sum_{n\geq0}e^{-\lambda_na}\) et une intégration sur les portions linéaires de \(\tau_\varepsilon\) donnent

\[
\int_0^\infty e^{-\lambda a}
\bigl[2\tau_\varepsilon(d)-\tau_\varepsilon(d-a)
-\tau_\varepsilon(d+a)\bigr]\,da
=\frac{2e^{-\lambda|d|}-e^{-\lambda|d-\varepsilon|}
-e^{-\lambda|d+\varepsilon|}}{\varepsilon\lambda^2}.
\]
Pour justifier l’échange de la série et de l’intégrale, la fonction triangulaire vérifie en outre

\[
|2\tau_\varepsilon(d)-\tau_\varepsilon(d-a)-\tau_\varepsilon(d+a)|
\leq\min(4,2a/\varepsilon).
\]
Son produit par \(\mu(a)\) est intégrable tant à l’origine qu’à l’infini. La somme des intégrales absolues est ainsi bornée, ce qui justifie l’interversion. Le terme intégré est également borné en valeur absolue par \(4/(\varepsilon\lambda_n^2)\), dont la série converge. L’intervention conjointe des deux côtés annule la singularité de \(\mu\) à l’origine. On peut donc sommer l’identité intégrée et obtenir le terme gamma déclaré. \hfill\(\square\)

Ces indicatrices sont des éléments du domaine fermé des lecteurs, et non des fonctions lisses. L’extension à partir de \(\mathcal D_R\) se justifie sans modifier la forme. Pour un intervalle normalisé,

\[
\|b_{\varepsilon,t}-T_ab_{\varepsilon,t}\|_2^2
=2\min(|a|/\varepsilon,1).
\]
L’intégrale contre \(\mu(a)\) est finie. Soit \(b_{\varepsilon,t}^{(\delta)}=\rho_\delta*b_{\varepsilon,t}\) une régularisation lisse, à support contenu dans une fenêtre fixe. On prend \(\rho\geq0\), lisse à support compact, avec \(\int\rho=1\), et \(\rho_\delta(x)=\delta^{-1}\rho(x/\delta)\). Le symbole gamma est

\[
\sigma_\Gamma(\xi)=\int_0^\infty
\mu(a)|1-e^{-ia\xi}|^2\,da.
\]
En Fourier, l'intégrande qui exprime la norme gamma de la différence contient \(\sigma_\Gamma|\widehat b|^2 |\widehat\rho(\delta\xi)-1|^2\). Le premier produit est intégrable et le dernier facteur est borné par \(4\) et tend vers zéro. La convergence dominée prouve la convergence gamma. Les canaux premiers sont finis et bornés dans cette fenêtre ; les canaux polaires sont des fonctionnelles continues sur son support. On obtient la convergence pour la norme conjointe des lecteurs et, par conséquent, celle de chaque matrice de Gram finie.

\HMTresultado{Corollaire}{Deux orientations d’une même fenêtre première}{rz-num-03-17}
Posons \(D_\varepsilon=K_\varepsilon(0)\). Pour deux centres séparés par \(L\), le Gram complet est

\[
G_2=\begin{pmatrix}D_\varepsilon&K_\varepsilon(L)\\
K_\varepsilon(L)&D_\varepsilon\end{pmatrix}.
\]
Par conséquent

\[
\mathscr W(af+bg,af+bg)=
\frac{D_\varepsilon+K_\varepsilon(L)}2|a+b|^2
+\frac{D_\varepsilon-K_\varepsilon(L)}2|a-b|^2.
\]
La positivité sur l’espace qu’elles engendrent exige simultanément \(D_\varepsilon+K_\varepsilon(L)\geq0\) et \(D_\varepsilon-K_\varepsilon(L)\geq0\).

Si \(L\geq\varepsilon\), on définit

\[
I_\varepsilon(L)=
\frac{F(L-\varepsilon)-2F(L)+F(L+\varepsilon)}{\varepsilon}>0,
\]
\[
P_\varepsilon(L)=
\sum_{|\ell_{p,k}-L|<\varepsilon}w_{p,k}
\left(1-\frac{|\ell_{p,k}-L|}{\varepsilon}\right).
\]
La convexité stricte de chaque exponentielle prouve la positivité de \(I_\varepsilon\). La condition des deux signes s’écrit exactement

\[
2a_\varepsilon^2\cosh(L/2)-I_\varepsilon(L)-D_\varepsilon
\leq P_\varepsilon(L)
\leq
2a_\varepsilon^2\cosh(L/2)-I_\varepsilon(L)+D_\varepsilon.
\]
Ainsi, le croisement complet contrôle une fenêtre de longueurs d’horloges irréductibles avec leurs poids, et non seulement une valeur diagonale. L’identité résulte de la proposition~\ref{rz-num-03-16} et de la diagonalisation de la matrice à deux entrées ; elle ne présuppose pas son signe. \hfill\(\square\)

\HMTresultado{Proposition}{Contrôle conjoint de neuf fonctions tests}{rz-num-03-18}
Soient \(\varepsilon=1/100\), \(t_j=j\log2\), \(0\leq j\leq8\), et \(G_{ij}=K_\varepsilon(t_i-t_j)\). Le calcul par intervalles conservé avec ce chapitre certifie

\[
\lambda_{\min}(G)>0.8479.
\]
En conséquence,

\[
\mathscr W\left(\sum_{j=0}^8z_jb_{\varepsilon,t_j},
\sum_{j=0}^8z_jb_{\varepsilon,t_j}\right)
>0.8479\sum_{j=0}^8|z_j|^2
\quad(z\neq0).
\]
\textbf{Démonstration et reproduction.} Pour \(F_N(x)=\sum_{n=0}^{N-1}e^{-\lambda_nx}/\lambda_n^2\), la fonction sommée décroît et

\[
0\leq F(x)-F_N(x)
\leq e^{-\lambda_Nx}
\left(\lambda_N^{-2}+\frac1{2\lambda_N}\right).
\]
En effet, on sépare le premier terme de la queue et on borne le reste par l’intégrale de \((2u+1/2)^{-2}\), en extrayant la plus grande exponentielle. On utilise \(N=1024\) et des opérations par intervalles avec arrondi vers l’extérieur. La condition \(8\log2+\varepsilon<\log259\) permet de recenser toutes les paires premier–puissance jusqu’à \(259\) : on obtient \(71\) lignes. Les irréductibles et les produits proviennent du producteur natif déjà conservé ; la reconnaissance par les entiers ordinaires est ultérieure. La constante \(c_\Gamma\) est évaluée à partir des intervalles certifiés des coordonnées précédentes, et non à partir d’une valeur cible de cette forme.

La formule de la proposition~\ref{rz-num-03-16} donne des intervalles pour les neuf entrées de la première ligne. Pour chaque ligne \(i\), on vérifie avec ces intervalles

\[
G_{ii}-\sum_{j\neq i}|G_{ij}|>0.8479.
\]
La matrice est hermitienne. L'inégalité \(2|z_i||z_j|\leq|z_i|^2+|z_j|^2\), appliquée à chaque terme croisé, donne directement la borne quadratique annoncée. La vérification conserve les intervalles des neuf lignes et le recensement complet. Les marges strictes et la convergence des matrices de Gram démontrée précédemment permettent également de régulariser simultanément les neuf fonctions tests, en obtenant la positivité sur l'espace lisse qu'elles engendrent pour une régularisation suffisamment fine. \hfill\(\square\)

Ce contrôle couvre toutes les combinaisons complexes de neuf fonctions tests, non uniquement leur somme. Sa dimension demeure finie. Le registre computationnel déclare cette portée et n’utilise pas son résultat comme certificat d’une inégalité universelle.

\HMTresultado{Proposition}{Raffinement et huit composantes de détail}{rz-num-03-19}
La subdivision géométrique d’un intervalle en neuf parties satisfait

\[
b_{\varepsilon,t}=\frac13\sum_{r=0}^8
b_{\varepsilon/9,t+(r-4)\varepsilon/9}.
\]
Si \(V=(1,\ldots,1)^{\mathsf T}/3\), alors

\[
G_{\mathrm{padre}}=V^*G_{\mathrm{hijos}}V.
\]
Pour plusieurs prédécesseurs, on prend l'inclusion isométrique par blocs avec cette même colonne. Dans une base formée de la composante uniforme et des huit composantes de détail par prédécesseur, écrivons

\[
G_{\mathrm{hijos}}=
\begin{pmatrix}A&B^*\\B&D\end{pmatrix},\qquad
A=G_{\mathrm{padres}}.
\]
Si \(A>0\), la condition exacte est

\[
G_{\mathrm{hijos}}\succeq0
\quad\Longleftrightarrow\quad D-BA^{-1}B^*\succeq0.
\]
\textbf{Démonstration.} Les neuf intervalles successeurs partitionnent le prédécesseur, à des extrémités de mesure nulle près. La normalisation de leurs indicateurs prouve l'identité linéaire, et la sesquilinéarité donne l'identité des matrices de Gram avec tous les termes croisés. Enfin,

\[
\begin{aligned}
\begin{pmatrix}x\\y\end{pmatrix}^*
\begin{pmatrix}A&B^*\\B&D\end{pmatrix}
\begin{pmatrix}x\\y\end{pmatrix}
={}&(x+A^{-1}B^*y)^*A(x+A^{-1}B^*y)\\
&+y^*(D-BA^{-1}B^*)y.
\end{aligned}
\]
Prendre \(x=-A^{-1}B^*y\) prouve la nécessité et la même identité prouve la suffisance. \hfill\(\square\)

La composante uniforme reproduit le Gram précédent ; les détails contiennent les nouvelles interférences. Conserver les deux parties est l'exigence exacte permettant de prolonger une preuve positive sous ce raffinement. Le retour de neuf phases ne transforme pas à lui seul le Gram en matrice circulante : dans la collection précédente, \(K_\varepsilon(\log2)\neq K_\varepsilon(8\log2)\). La phase et la mémoire des distances restent des données distinctes.

La matrice de détail et son complément de Schur fournissent une condition vérifiable pour chaque extension. Le lecteur prospectif n'est pas défini au moyen d'une factorisation rétrospective de la forme ; on vérifie la forme exprimée par les lecteurs précédemment construits. Pour compléter l'argument global, une estimation sur tous ces détails est requise, uniforme dans les extensions cofinales et sur le domaine complet des fonctions tests.

\subsection{Assemblage coercif et stabilité des détails nonadiques}
\label{rz:ensamblaje}
La propriété basale permet de contrôler des espaces complets de fonctions. Pour assembler plusieurs intervalles, il faut comparer cette marge locale à chaque interaction entre routes. Le noyau complet de la section précédente fournit une estimation de ces interactions avant de choisir une fonction. On obtient ainsi un résultat de dimension infinie sur un domaine non connexe, non seulement un gramien d’indicatrices.

Les positions sont prises après la représentation des horloges irréductibles. Le cas explicite utilise les translations successives de l'horloge native 2 et une largeur positionnelle déclarée. Les opérateurs de différence conservent les deux orientations. Le raffinement ultérieur subdivise chaque intervalle en neuf parties et maintient la même fonction avant ses deux lectures d'énergie. Cette collection analytique est spécifiée comme telle : son cardinal ne l'identifie pas à lui seul à une sélection canonique d'états TPK.

\HMTresultado{Proposition}{Marge diagonale et borne des termes croisés}{rz-num-03-20}
Soient \(t_1<\cdots<t_m\) des centres distincts et \(I_j=(t_j-\varepsilon/2,t_j+\varepsilon/2)\) des intervalles disjoints, avec \(0<\varepsilon<\log2\). Définissons

\[
\delta_\varepsilon=2M_\Gamma(\varepsilon)+c_\Gamma
-2\sinh(\varepsilon/2)+\varepsilon.
\]
Pour \(f_j\in C_c^\infty(I_j)\),

\[
\mathscr W(f_j,f_j)\geq\delta_\varepsilon\|f_j\|_2^2.
\]
Si \(d_{ij}=|t_i-t_j|>\varepsilon\), posons

\[
\rho_{ij}(\varepsilon)=
\sum_{|\ell_{p,k}-d_{ij}|<\varepsilon}w_{p,k}
+\varepsilon\mu(d_{ij}-\varepsilon)
+2\varepsilon\cosh((d_{ij}+\varepsilon)/2).
\]
Alors

\[
|\mathscr W(f_i,f_j)|
\leq\rho_{ij}(\varepsilon)\|f_i\|_2\|f_j\|_2
\qquad(i\neq j).
\]
\textbf{Démonstration.} Dans un intervalle de largeur \(\varepsilon\), la queue gamma apporte au moins \(2M_\Gamma(\varepsilon)\|f_j\|_2^2\). Il n'y a pas de corrélation premier–puissance diagonale, car sa longueur minimale est \(\log2>\varepsilon\). Dans une fenêtre ambiante plus grande, chaque différence première et son contre-terme s'annulent sur ce support. Après centrage de l'intervalle, le port polaire négatif a pour carré de sa norme \(2\sinh(\varepsilon/2)-\varepsilon\). Écarter le port polaire positif et la partie gamma précédant la queue donne la borne. L'invariance par translation de la forme complète transporte la même borne à n'importe quel centre.

Dans le croisement, la contribution scalaire est nulle par disjonction des supports. Pour chaque longueur positive, une seule des deux orientations peut relier les deux intervalles. L’unitarité de la translation borne la contribution première par la somme des poids indiquée. Le croisement gamma a pour noyau \(-\mu(|x-y|)\) sur \(I_i\times I_j\). Comme \(\mu\) est décroissante,

\[
|\mathscr W_\Gamma(f_i,f_j)|
\leq\mu(d_{ij}-\varepsilon)\|f_i\|_1\|f_j\|_1
\leq\varepsilon\mu(d_{ij}-\varepsilon)\|f_i\|_2\|f_j\|_2.
\]
Le noyau polaire complet est \(2\cosh((x-y)/2)\). La borne \(|x-y|\leq d_{ij}+\varepsilon\) et la même inégalité entre les normes \(L^1\) et \(L^2\) donnent le terme polaire. Additionner les trois estimations prouve le résultat. \hfill\(\square\)

\HMTresultado{Théorème}{Coercivité de l’assemblage}{rz-num-03-21}
Si

\[
\eta=\delta_\varepsilon-
\max_i\sum_{j\neq i}\rho_{ij}(\varepsilon)>0,
\]
alors, pour \(\Omega=\bigcup_j I_j\),

\[
\boxed{\mathscr W(f,f)\geq\eta\|f\|_2^2
\quad\text{pour toute }f\in C_c^\infty(\Omega).}
\]
L’affirmation vaut dans la fermeture de ce domaine pour la norme conjointe des lecteurs. En particulier, pour

\[
m=9,\qquad t_j=j\log2\ (0\leq j\leq8),\qquad
\varepsilon=1/1000,
\]
l’évaluation certifiée donne \(\eta>1\).

\textbf{Démonstration.} En écrivant \(f=\sum_jf_j\) sur les supports disjoints et \(x_j=\|f_j\|_2\), on obtient

\[
\begin{aligned}
\mathscr W(f,f)&\geq
\delta_\varepsilon\sum_jx_j^2
-2\sum_{i<j}\rho_{ij}(\varepsilon)x_ix_j\\
&\geq\sum_i\left(\delta_\varepsilon-
\sum_{j\neq i}\rho_{ij}(\varepsilon)\right)x_i^2.
\end{aligned}
\]
On a utilisé \(2x_ix_j\leq x_i^2+x_j^2\) pour chaque croisement. La somme des carrés des normes est \(\|f\|_2^2\). La continuité des lecteurs étend l’inégalité au domaine fermé.

Dans le cas explicite, les distances sont \(d\log2\), \(1\leq d\leq8\). Pour \(\varepsilon=1/1000\),

\[
2^d(e^\varepsilon-1)<\frac{256}{999}<1,
\qquad 2^d(1-e^{-\varepsilon})<\frac{256}{1000}<1.
\]
La première inégalité utilise la série exponentielle et \(e^\varepsilon<1/(1-\varepsilon)\) ; la seconde utilise \(1-e^{-\varepsilon}<\varepsilon\). Les puissances exprimées étant entières, la seule horloge de puissance première dans la fenêtre de cette distance est \((p,k)=(2,d)\). Alors

\[
\rho_d=(\log2)2^{-d/2}
+\frac1{1000}\mu(d\log2-1/1000)
+\frac2{1000}\cosh((d\log2+1/1000)/2).
\]
Le producteur natif vérifie en outre le recensement complet jusqu'à \(259\). Le calcul de \(\delta_\varepsilon\) utilise la formule déjà démontrée de \(M_\Gamma\) et l'identité

\[
\arctan t=\frac\pi4-
\arctan\frac{1-t}{1+t}\qquad(0<t<1).
\]
Le second argument est petit et sa série alternée admet pour borne le premier terme omis. La représentation interne de \(\pi\), l'intervalle rationnel de gamma et les opérations ultérieures sur les intervalles produisent les bornes

\[
\delta_\varepsilon>4.4921,\qquad
\max_i\sum_{j\neq i}\rho_{ij}<2.5385,\qquad
\eta>1.9536>1.
\]
Le reçu conserve la marge de chaque ligne. On n'évalue pas de fonctions particulières pour en inférer l'inégalité universelle sur \(C_c^\infty(\Omega)\) : celle-ci découle de l'estimation opératorielle précédente. \hfill\(\square\)

\HMTresultado{Corollaire}{Tous les raffinements du domaine assemblé}{rz-num-03-22}
Subdivisons chaque \(I_j\) successivement en neuf parties. Les indicateurs normalisés des \(9^{n+1}\) intervalles de profondeur \(n\) forment une collection orthonormale dans \(L^2\). Leur Gram complet satisfait

\[
G_n\succeq\eta I\qquad(n\geq0).
\]
L'inclusion uniforme \(V_n\), avec neuf composantes \(1/3\) par prédécesseur, vérifie

\[
V_n^*G_{n+1}V_n=G_n.
\]
Dans la décomposition uniforme–détail,

\[
G_{n+1}=\begin{pmatrix}A_n&B_n^*\\B_n&D_n\end{pmatrix},
\qquad A_n=G_n,
\]
on obtient, à toute profondeur,

\[
\boxed{D_n-B_nA_n^{-1}B_n^*\succeq\eta I.}
\]
\textbf{Démonstration.} Le théorème s’applique aux régularisations de chaque combinaison d’indicatrices. La convergence en norme des lecteurs de la sous-section~\ref{rz:ventanas} permet de passer à la limite. Une indicatrice qui atteint une extrémité de \(I_j\) est d’abord approchée en contractant légèrement son support vers le centre, puis en régularisant à l’intérieur de \(I_j\). La corrélation triangulaire et sa borne intégrable de translation prouvent la même convergence ; on n’utilise pas d’extension du domaine de coercivité. L’orthonormalité dans \(L^2\) transforme la borne fonctionnelle en \(G_n\succeq\eta I\), et la subdivision exacte donne l’identité de compression.

Pour un vecteur de détail \(y\), on prend \(x=-A_n^{-1}B_n^*y\). L'identité de Schur donne

\[
\begin{aligned}
y^*(D_n-B_nA_n^{-1}B_n^*)y
&=\begin{pmatrix}x\\y\end{pmatrix}^*G_{n+1}
\begin{pmatrix}x\\y\end{pmatrix}\\
&\geq\eta(\|x\|^2+\|y\|^2)\geq\eta\|y\|^2.
\end{aligned}
\]
Le même \(\eta\) convient à toute profondeur. \hfill\(\square\)

Le corollaire contrôle les huit nouveaux détails par prédécesseur à profondeur arbitraire dans l'assemblage. La fonction n'est pas remplacée par sa moyenne et ses composantes de somme nulle ne sont pas écartées. C'est un renforcement du Gram de neuf indicateurs : il couvre des espaces fonctionnels de dimension infinie et donne une borne indépendante de la profondeur.

\subsubsection*{Collections finies et extension du domaine}
Pour tout ensemble fini de centres distincts, il existe une largeur positive suffisamment petite permettant d'appliquer le théorème~\ref{rz-num-03-21}. En effet, \(2M_\Gamma(\varepsilon)\to+\infty\) lorsque \(\varepsilon\downarrow0\), tandis que le terme polaire négatif local tend vers zéro. Les distances entre centres sont séparées de zéro. Les sommes de poids premiers restent bornées lorsque chaque fenêtre se réduit, car il existe un nombre fini de longueurs en dessous de toute borne fixée. Les termes croisés gamma et polaire tendent vers zéro. Par conséquent, \(\delta_\varepsilon\) finit par dépasser la somme des interactions de chaque ligne.

La largeur admissible dépend des centres. Les unions d’intervalles étroits obtenues de cette manière ne forment pas, par la seule construction précédente, un recouvrement cofinal de toutes les fonctions à support compact. La positivité globale exige de contrôler également la croissance du domaine et ses nouvelles interactions. L’uniformité démontrée en profondeur et l’extension à tout domaine de fonctions tests, non démontrée ici, restent distinguées.

\subsection{Coercivité des détails dans une fenêtre arbitraire fixée}
\label{rz:detalles}
L’estimation précédente conserve une marge uniforme lors du raffinement de l’intérieur de neuf intervalles séparés. Une seconde construction permet de contrôler les détails fins de toute fenêtre fixée, y compris ses cellules contiguës. Les deux affirmations ont des domaines différents et complémentaires : la première estime la fonction complète sur un support non connexe ; la seconde estime les composantes de moyenne nulle par cellule dans une fenêtre complète.

Les horloges, poids et lecteurs construits précédemment sont maintenus. En particulier, les longueurs \(\ell_{p,k}\), les poids \(w_{p,k}\), le scalaire \(c_\Gamma\) et la représentation interne \(\pi\) conservent leur provenance. L'argument suivant agit sur la réalisation fonctionnelle de ces objets ; il ne modifie pas leur générateur. Écrivons

\[
\beta_R=2\sinh(R/2)-R,\qquad
\mathfrak c_R=\frac23-c_\Gamma+
2\sum_{\ell_{p,k}\leq R}w_{p,k}+\beta_R.
\]
La somme est finie et provient du produit natif sous la borne \(e^R\). Divisons \(I_R=(-R/2,R/2)\) en \(M\) cellules de longueur au plus \(h\), avec \(0<h\leq1\).

\HMTresultado{Théorème}{Seuil explicite pour les détails de moyenne nulle}{rz-num-03-23}
Soit \(d\) une fonction du domaine conjoint des lecteurs, à support dans \(I_R\), dont l'intégrale dans chaque cellule est nulle. Alors

\[
\boxed{\mathscr W_R(d,d)\geq
\bigl[\log(1+4/h)-\mathfrak c_R\bigr]\|d\|_2^2.}
\]
En conséquence, la forme est strictement coercive dans cet espace de détails si

\[
0<h<\min\left(1,\frac4{\exp(\mathfrak c_R)-1}\right).
\]
Le seuil dépend de la fenêtre \(R\), mais non de la fonction ni d’une énumération finie d’exemples.

\textbf{Démonstration.} Dans chaque cellule \([a,b]\), on définit \(F(x)=\int_a^x d(t)\,dt\). La condition de moyenne nulle implique \(F(a)=F(b)=0\). En réunissant les primitives, on obtient \(F\in H_0^1(I_R)\), prolongée par zéro hors de la fenêtre, avec \(F'=d\). L’inégalité de Dirichlet dans chaque cellule donne

\[
\|F\|_2^2\leq\frac{h^2}{\pi^2}\|d\|_2^2.
\]
Ici, \(\pi\) est la même coordonnée interne reconnue dans la réalisation fonctionnelle. Aucune nouvelle constante n'est sélectionnée.

La queue gamma admet le développement positif déjà utilisé

\[
\mu(a)=\sum_{n=0}^{\infty}e^{-\lambda_n a},
\qquad \lambda_n=2n+\frac12.
\]
Pour \(\lambda>0\), soit \(C_\lambda\) la convolution par \(e^{-\lambda|x|}\). La contribution d’un terme est

\[
E_\lambda(d)=\int_0^\infty e^{-\lambda a}
\|(I-T_a)d\|_2^2\,da
=\frac2\lambda\|d\|_2^2-
\langle d,C_\lambda d\rangle\geq0.
\]
Sa positivité est celle d'une intégrale de carrés de normes. Notons \(\mathcal F_u F=(2\pi)^{-1/2}\widehat F\) la transformée de Fourier unitaire, en la distinguant de la transformée non normalisée employée précédemment. Le multiplicateur de \(C_\lambda\) est \(2\lambda/(\lambda^2+\xi^2)\). Comme \(d=F'\), Plancherel implique

\[
\begin{aligned}
\langle d,C_\lambda d\rangle
&=\int_{\mathbb R}
\frac{2\lambda\xi^2}{\lambda^2+\xi^2}
|(\mathcal F_u F)(\xi)|^2\,d\xi\\
&\leq2\lambda\|F\|_2^2
\leq\frac{2\lambda h^2}{\pi^2}\|d\|_2^2.
\end{aligned}
\]
Par Tonelli, la forme gamma est la somme des \(E_{\lambda_n}\). On conserve les \(N+1\) premiers termes et on écarte uniquement une queue non négative. Il en résulte

\[
\begin{aligned}
Q_\Gamma(d)&\geq S_N(h)\|d\|_2^2,\\
S_N(h)&=4\sum_{n=0}^N\frac1{4n+1}
-\frac{h^2(N+1)(2N+1)}{\pi^2}.
\end{aligned}
\]
Choisissons \(N=\lfloor1/h\rfloor\). Alors \(\lambda_N h\leq2+h/2<\pi\) ; les bornes retenues sont non négatives. La comparaison avec l’intégrale d’une fonction décroissante donne

\[
4\sum_{n=0}^N\frac1{4n+1}
\geq\int_0^{N+1}\frac4{4x+1}\,dx
=\log(4N+5)\geq\log(1+4/h).
\]
De plus, \(\pi^2>9\) et \(h\leq1\) impliquent

\[
\frac{h^2(N+1)(2N+1)}{\pi^2}
\leq\frac{(1+h)(2+h)}{\pi^2}<\frac23.
\]
Par conséquent, \(Q_\Gamma(d)\geq[\log(1+4/h)-2/3]\|d\|_2^2\). Dans la forme complète, le terme scalaire apporte \(c_\Gamma\|d\|_2^2\) ; la partie première est minorée par \(-2\sum_{\ell_{p,k}\leq R}w_{p,k}\|d\|_2^2\) au moyen de l’unitarité de chaque translation ; et la partie polaire est bornée par \(-\beta_R\|d\|_2^2\), comme démontré dans la sous-section~\ref{rz:gamma}. Leur somme prouve l’inégalité annoncée. Le seuil écrit équivaut à \(\log(1+4/h)>\mathfrak c_R\). Le passage au domaine fermé conserve l’estimation par continuité des lecteurs et des intégrales de cellule. \hfill\(\square\)

\HMTresultado{Corollaire}{Dimension finie des directions négatives possibles}{rz-num-03-24}
Étant fixés \(R\) et une partition satisfaisant le seuil précédent, tout sous-espace du domaine sur lequel \(\mathscr W_R\) est strictement définie négative a une dimension au plus égale à \(M\), le nombre de cellules.

\textbf{Démonstration.} Considérons l'application linéaire

\[
\mathcal P(d)=\left(\int_{I_1}d,\ldots,
\int_{I_M}d\right)\in\mathbb C^M.
\]
Son noyau est l'espace des détails du théorème~\ref{rz-num-03-23}, sur lequel la forme est définie positive. La restriction de \(\mathcal P\) à un sous-espace strictement négatif est injective : un vecteur non nul de son noyau aurait à la fois une valeur positive et négative. La dimension de ce sous-espace ne peut dépasser \(M\). \hfill\(\square\)

Cette réduction conserve une question finie de couplage dans chaque fenêtre. Elle ne revient pas à prouver que l’indice négatif est nul. Dans une décomposition moyenne–détail, le bloc des détails est coercif à une profondeur suffisante, mais il faut conserver le bloc des moyennes et les termes croisés. Si l’on utilise le complément de Schur, son signe exige cette composition complète. En particulier, une résonance première exacte entre deux cellules ne disparaît pas en imposant une moyenne nulle : elle peut agir comme un multiple négatif de l’identité sur leurs détails correspondants. Le théorème maintient de tels termes grâce à sa borne de translation.

Le résultat apporte une seconde uniformité effective : dans toute fenêtre fixée, les détails suffisamment fins ont une borne commune explicite. La sous-section~\ref{rz:ensamblaje} prouve en outre le signe de l'assemblage complet, avec ses termes croisés, pour la réunion des neuf intervalles spécifiée. L'extension de ce signe à toute fenêtre reste l'étape globale que cette édition ne tient pas pour démontrée.
